\documentclass[11pt,a4paper]{article}

% --- Paquetes de Diseño de Página y Márgenes ---
\usepackage[margin=1in]{geometry}
\usepackage{authblk}

% --- Paquetes Matemáticos, Gráficos y Utilidades ---
\usepackage{amsmath, amssymb, amsthm}
\usepackage{mathtools}
\usepackage{bm}
\usepackage{dsfont}
\usepackage{braket}
\usepackage{physics}
\usepackage{graphicx}
\usepackage{color}
\usepackage{xcolor}
\usepackage{booktabs}
\usepackage{scrhack}

% --- Algoritmos y Código ---
\usepackage{algorithm}
\usepackage{algpseudocode}
\usepackage{listings}

\definecolor{codegreen}{rgb}{0,0.6,0}
\definecolor{codegray}{rgb}{0.5,0.5,0.5}
\definecolor{codepurple}{rgb}{0.58,0,0.82}
\definecolor{backcolour}{rgb}{0.95,0.95,0.92}

\lstdefinestyle{mystyle}{
    backgroundcolor=\color{backcolour},   
    commentstyle=\color{codegreen},
    keywordstyle=\color{magenta},
    numberstyle=\tiny\color{codegray},
    stringstyle=\color{codepurple},
    basicstyle=\ttfamily\footnotesize,
    breakatwhitespace=false,         
    breaklines=true,                 
    captionpos=b,                    
    keepspaces=true,                 
    numbers=left,                    
    numbersep=5pt,                  
    showspaces=false,                
    showstringspaces=false,
    showtabs=false,                  
    tabsize=2,
    language=Python
}
\lstset{style=mystyle}

% --- Hipervínculos ---
\usepackage[colorlinks=true, linkcolor=blue!70!black, citecolor=red!70!black, urlcolor=blue!70!black]{hyperref}
\usepackage{orcidlink}
\newcommand{\orcid}[1]{\orcidlink{#1}}

% --- Cleveref ---
\usepackage{cleveref}
\crefname{equation}{Eq.}{Eqs.}
\crefname{section}{Sec.}{Secs.}
\crefname{table}{Table}{Tables}
\crefname{figure}{Fig.}{Figs.}

% --- Entornos de teoremas ---
\newtheorem{theorem}{Theorem}[section]
\newtheorem{lemma}[theorem]{Lemma}
\newtheorem{corollary}[theorem]{Corollary}
\newtheorem{definition}[theorem]{Definition}
\newtheorem{remark}[theorem]{Remark}

% --- Macros personalizados ---
\newcommand{\Z}{\mathbb{Z}}
\newcommand{\C}{\mathbb{C}}
\newcommand{\R}{\mathbb{R}}
\newcommand{\N}{\mathbb{N}}
\newcommand{\Hilb}{\mathcal{H}}
\newcommand{\Zmod}{\mathbb{Z}/6\mathbb{Z}}
\DeclareMathOperator{\IPR}{IPR}

\title{Multifractal non-ergodic extended phase in power-law random banded matrices with modular arithmetic constraints}

\author{José Ignacio Peinador Sala \orcidlink{0009-0008-1822-3452}}
\affil{Independent Researcher, Valladolid, Spain}
\affil{\text{Email: }\texttt{joseignacio.peinador@gmail.com}}

\date{\today}

\begin{document}

\maketitle

We introduce a deterministic-stochastic Hamiltonian for a one-dimensional quantum lattice whose hoppings are constrained by coprimality with respect to 6, naturally encoded in $\mathbb{Z}/6\mathbb{Z}$. The diagonal potential implements the exact inversion of the Riemann--von Mangoldt counting function via the Lambert $W$ function, while the off-diagonal disorder belongs to the power-law random banded matrix (PRBM) class with decay exponent $\nu=0.75$. Exact diagonalization up to $N=20\,000$ shows that the macroscopic density follows the Weyl law and local fluctuations obey GUE statistics. However, the long-range spectral rigidity reveals a sub-diffusive spectral form factor ramp ($\gamma \approx 0.609$) and a suppressed generalized fractal dimension ($D_2 \approx 0.243$), placing the system deep inside a non-ergodic extended (NEE) phase driven by the arithmetic selection rule. This mask provides topological protection against localization and enforces a rigorous upper bound $D_2 \le \varphi(m)/m$, which we mechanically certify using the Lean 4 theorem prover. Massive GPU diagonalization confirms that individual eigenstates sharply respect the refined bound $D_2 \approx 1/4$. In contrast, actual Riemann zeros are ergodic ($\gamma \approx 1.12$), confirming that the observed multifractality is an intrinsic property of the modular constraint.

\section{Introduction}

Random matrix theory (RMT) and the physics of disordered quantum systems have long been intertwined with number theory, most famously through the statistical properties of the Riemann zeta zeros. Montgomery's pair correlation conjecture and Odlyzko's numerical computations demonstrated that the zeros on the critical line exhibit the same level repulsion as eigenvalues of Gaussian unitary ensemble (GUE) matrices \cite{Montgomery1973,Odlyzko2001}. This motivated the Hilbert--P\'olya idea that the imaginary parts of the zeros might be eigenvalues of a Hermitian operator, a program that remains open.

Rather than attempting to construct such an operator, we take a different route: we use the smooth density of zeros — the Riemann--von Mangoldt formula — as the diagonal potential of a one-dimensional lattice Hamiltonian, and then add off-diagonal disorder subject to an arithmetic constraint derived from the ring $\mathbb{Z}/6\mathbb{Z}$. The resulting model is a deterministic-stochastic hybrid whose thermodynamic and dynamical properties can be studied in their own right.

The main novelty is the arithmetic selection rule: hopping between sites $i$ and $j$ is allowed only if the distance $|i-j|$ is coprime to~6, i.e., belongs to the residue classes $1$ or $5$ modulo~6. This rule eliminates exactly two thirds of the connections and introduces a periodic sparsification of the interaction graph. The system falls into the universality class of power-law random banded matrices (PRBM) with spatial decay $\nu=0.75$, a value that places it at the center of the extended chaotic phase where eigenstates are neither fully ergodic nor localized \cite{Mirlin1996}.

We show by large-scale exact diagonalization that the macroscopic spectral density remains locked to the input Weyl law, while the local level statistics are GUE. However, the long-range spectral rigidity, quantified by the spectral form factor (SFF) and the generalized fractal dimension $D_2$, deviates strongly from the ergodic limit. The SFF ramp exponent is $\gamma \approx 0.609$, far below the standard GUE value $\gamma=1$, and the fractal dimension is $D_2 \approx 0.243$, indicating that wavefunctions populate a sparse multifractal support. These two independent metrics consistently place the system in a non-ergodic extended (NEE) phase.

We further demonstrate that $D_2$ and $\gamma$ are systematically suppressed as the density of surviving hopping channels $\varphi(m)/m$ decreases. This channel-density scaling, together with a rigorous upper bound $D_2 \le \varphi(m)/m$, establishes a direct connection between the arithmetic restriction and the emergent multifractality. To ensure absolute mathematical rigor, the structural topology of the interaction graph and the algebraic limits governing this phase have been formally verified using the Lean~4 proof assistant. We also show that the arithmetic mask provides topological protection against Anderson localization, preserving GUE spectral rigidity well beyond the unconstrained critical point.

To test whether this sub-diffusive behavior might reflect an unknown property of the Riemann zeros themselves, we compute the SFF of the first $10^4$ actual zeros. We find $\gamma \approx 1.12$, completely consistent with ordinary GUE. The discrepancy confirms that the observed NEE phase is an intrinsic feature of the modular constraint and the PRBM structure, not a spectral fingerprint of the zeta function.

The model presented here is therefore a transparent example of how a simple arithmetic filter can generate multifractal non-ergodic behavior in a quantum system. The $\mathbb{Z}/6\mathbb{Z}$ mask is not an arbitrary choice; it emerges naturally as the center of the Standard Model gauge group in noncommutative geometry, where the internal space has KO-dimension 6 modulo 8 \cite{Connes2008,Chamseddine2007}. This geometric origin provides a deeper motivation for the arithmetic constraint.

The paper is organized as follows. Section~\ref{sec:model} defines the Hamiltonian and its parameters. Section~\ref{sec:results} presents the numerical results, including the macroscopic alignment, the GUE level statistics, the SFF and $D_2$ analysis, the channel-density scaling, the topological protection, and the comparison with the actual Riemann zeros. Section~\ref{sec:discussion} discusses the physical interpretation and the connection to noncommutative geometry. Technical derivations are collected in the Appendix, including the formal bound $D_2 \le \varphi(m)/m$.

\section{Model Hamiltonian}
\label{sec:model}

\subsection{Hilbert space and diagonal potential}

We consider a quantum system on the Hilbert space $\Hilb = \ell^2(\mathbb{N})$ with orthonormal basis $\{|k\rangle\}_{k\ge 1}$. The unperturbed Hamiltonian is diagonal,
\begin{equation}
\hat H_0 |k\rangle = E_k |k\rangle,
\end{equation}
with the energies $E_k$ chosen to reproduce the macroscopic density of a generic sequence with a logarithmic level density, such as the Riemann zeros. The Riemann--von Mangoldt formula states that the number of zeros with imaginary part less than $E$ is
\begin{equation}
N(E) = \frac{E}{2\pi}\log\left(\frac{E}{2\pi e}\right) + \frac{7}{8} + O(E^{-1}).
\end{equation}
Setting $N(E_k)=k$ and inverting exactly by the Lambert $W$ function gives the deterministic diagonal term
\begin{equation}
\boxed{ E_k = \frac{2\pi (k - 7/8)}{W\!\left(\frac{k-7/8}{e}\right)} }.
\label{eq:diag}
\end{equation}
The constant $7/8$ is the Maslov phase originating from the Riemann--Siegel theta function \cite{Corless1996}. In the unfolded spectrum the mean level spacing tends to $2\pi$.

\subsection{Off-diagonal disorder and modular constraint}

The perturbation $\hat V$ is a random matrix with elements
\begin{equation}
V_{ij} = \epsilon \, d^{-\nu} \, G_{ij} \, \Xi(d),
\qquad d = |i-j|,
\end{equation}
where $G_{ij}$ are independent complex Gaussian variables with $\mathbb{E}[|G_{ij}|^2]=1$, $\epsilon$ is the coupling constant, $\nu wholesalers=0.75$ the spatial decay exponent, and $\Xi(d)$ is the arithmetic mask
\begin{equation}
\Xi(d) = \begin{cases}
1 & \text{if } \gcd(d,6)=1,\\
0 & \text{otherwise}.
\end{cases}
\end{equation}
This mask restricts hopping to distances $d\equiv1,5\pmod6$, thereby blocking all connections whose length shares a factor with~6. The total Hamiltonian
\begin{equation}
\hat H = \hat H_0 + \hat V
\label{eq:fullH}
\end{equation}
is manifestly Hermitian because $V_{ji} = \overline{V_{ij}}$.

\subsection{Parameters: \texorpdfstring{$\nu$}{nu}, \texorpdfstring{$\epsilon$}{epsilon}, and self-adjointness}

The operator $\hat V$ belongs to the PRBM class. The Kato--Rellich theorem requires that $\hat V$ be relatively bounded with respect to $\hat H_0$. Since $\sum_d d^{-2\nu}$ converges only for $\nu>0.5$, we must have $\nu>0.5$. To remain well inside the extended chaotic phase and avoid both thermal divergence ($\nu\le0.5$) and Anderson localization ($\nu\ge1$), we fix
\begin{equation}
\nu = 0.75,
\end{equation}
the exact center of the critical window $0.5<\nu<1$ \cite{Mirlin1996}. This value guarantees essential self-adjointness on the dense domain $\mathcal{D}(\hat H_0)$.

The coupling $\epsilon$ is not a free parameter but is determined by a local resonance condition. In the unfolded spectrum the mean spacing is $D=2\pi$. Two neighboring levels hybridize effectively when the variance of the coupling compensates the level spacing. Because the interaction is Hermitian, a transition $|n\rangle\to|m\rangle$ has a conjugate return path, doubling the effective variance in the local $2\times2$ subspace. Hence we impose
\begin{equation}
2\,\mathbb{E}[|V_{n,n+1}|^2] = D^2 \;\Longrightarrow\; 2\epsilon^2 = (2\pi)^2,
\end{equation}
yielding
\begin{equation}
\boxed{\epsilon_c = \pi\sqrt{2} \approx 4.44288}.
\label{eq:epsilon}
\end{equation}
This value lies at the threshold of maximal chaos in the local subspace \cite{Cunden2018,Gross1980}.

\subsection{Connection to noncommutative geometry}

Although the model is mathematically self-contained, the choice of modulus~6 has a deep physical motivation. In Connes' noncommutative geometry approach to the Standard Model, the finite internal space is described by the algebra $\mathcal{A}_F = \mathbb{C}\oplus\mathbb{H}\oplus M_3(\mathbb{C})$, whose unimodular gauge group is $G_{SM} = U(1)\times SU(2)\times SU(3)/\mu_6$. The center of this group is exactly $\mathbb{Z}/6\mathbb{Z}$ \cite{Connes2008}. Moreover, to avoid fermion doubling and produce the correct chiral spectrum, the internal space must have KO-dimension $6$ modulo $8$ \cite{Chamseddine2007}. Our arithmetic mask, which selects the two residues coprime to~6, can be seen as a discrete realization of the chiral grading of the internal space. The model thus inherits a geometric rationale that goes beyond algorithmic sieving.

\subsection{Discrete Dirac operator and the \texorpdfstring{$\mathbb{Z}/6\mathbb{Z}$}{Z/6Z} superselection rule from noncommutative geometry}
\label{sec:ncg_discrete}

The connection between the modular ring $\Zmod$ and the finite geometry of the Standard Model can be made quantitatively precise by considering the discretization of the Dirac operator on a one‑dimensional lattice that inherits the algebraic structure of the internal space $F$. We outline here the essential steps that lead to the arithmetic mask $\Xi(d)$ and to the qualitative form of the diagonal potential $E_k$.

In the spectral triple $(\mathcal{A}, \mathcal{H}, D)$ describing the full Standard Model \cite{Connes2008,Chamseddine2007}, the algebra is $\mathcal{A}=C^\infty(M)\otimes \mathcal{A}_F$ with $\mathcal{A}_F=\mathbb{C}\oplus\mathbb{H}\oplus M_3(\mathbb{C})$, and the Dirac operator decomposes as $D = D_M\otimes 1 + \gamma_M\otimes D_F$. The finite part $D_F$ acts on the fermionic Hilbert space $\mathcal{H}_F$ and encodes the Yukawa couplings and the CKM matrix.

The unimodular gauge group is $G_{SM}=U(1)\times SU(2)\times SU(3)\,/\,\mu_6$, where the discrete group $\mu_6\simeq\Zmod$ acts by multiplication with sixth roots of unity. This quotient is necessary to cancel the common center of the three factors and to obtain the correct hypercharge assignments. All physical fields must be invariant under this central action, which imposes a superselection rule: any transition between two states whose $U(1)$ charges differ by an integer that is not a multiple of~6 is forbidden.

To translate this restriction into a one‑dimensional lattice model, we consider a simplified setting where the continuous spacetime $M$ is replaced by a discrete circle $\mathbb{Z}_L$ with $L\to\infty$, and the internal space is projected onto a single direction that carries the charge information. Let $\psi_n$ be a fermionic amplitude localized at site $n\in\mathbb{Z}$. The Dirac operator on the lattice can be written schematically as
\begin{equation}
(D\psi)_n = \sum_{d\neq 0} \frac{\gamma(d)}{d^{\alpha}} \, \psi_{n+d} + \text{diagonal terms},
\end{equation}
where $\gamma(d)$ are matrices in the internal space and the power‑law decay $d^{-\alpha}$ mimics the long‑range nature of the Euclidean propagator on a one‑dimensional manifold (see Ref.~\cite{Mirlin1996} for related PRBM constructions).

The invariance under the central $\mu_6$ forces the transition amplitude between sites $n$ and $n+d$ to carry a phase factor that cancels exactly the action of $\mu_6$ on the internal indices. In the basis where the $U(1)$ generator is diagonal, the condition becomes
\begin{equation}
e^{2\pi i q d /6} = 1 \quad \text{for all allowed transitions},
\end{equation}
where $q$ is the $U(1)$ charge of the fermion. This equation requires $\gcd(d,6)=1$, because only then can the phase be absorbed by a suitable choice of the internal matrix $\gamma(d)$. Hence the off‑diagonal entries are restricted to distances $d$ that are coprime to~6. The most general charge‑invariant kinetic term therefore reads
\begin{equation}
(D\psi)_n = \sum_{\substack{d\neq 0 \\ \gcd(d,6)=1}} \frac{\Gamma_d}{d^{\nu}} \, \psi_{n+d} + \text{h.c.},
\end{equation}
where $\Gamma_d$ are now complex random matrices (GUE) that encode the remaining geometric information of the finite space $F$. Identifying the coupling $\epsilon$ with the overall scale of $D_F$ and $\nu$ with the fall‑off exponent of the propagator, we recover exactly the off‑diagonal part of Eq.~\eqref{eq:fullH}.

The diagonal potential $E_k$ in Eq.~\eqref{eq:diag} can similarly be motivated from the NCG framework. The finite Dirac operator $D_F$ has a block‑diagonal form whose eigenvalues are related to the fermion masses. After integrating out the continuous degrees of freedom and taking the logarithmic large‑$N$ limit, the effective eigenvalue density becomes proportional to the Riemann‑von Mangoldt formula \cite{Connes1999}. The exact Lambert $W$ inversion used here is a discrete analogue of this asymptotic density, with the $7/8$ shift playing the role of the Maslov index that corrects the boundary terms in the trace formula.

Therefore, the full Hamiltonian $\hat H = \hat H_0 + \hat V$ can be regarded as a finite‑dimensional approximation of the Dirac operator $D$ projected onto a one‑dimensional lattice that retains the essential $\mathbb{Z}/6\mathbb{Z}$ central invariance. The arithmetic mask $\Xi(d)=\chi_0^{(6)}(d)$ is thus not an algorithmic artefact but the direct imprint of the KO-dimension 6 chiral grading on the discrete kinetic term. This geometric origin provides a physical motivation for why the model exhibits robust non-ergodic extended behavior: the central symmetry suppresses a large fraction of the hopping channels, naturally driving the system into the multifractal NEE phase characterized by $D_2\approx\varphi(6)/6=1/3$.

\section{Numerical results}
\label{sec:results}

\subsection{Exact diagonalization and macroscopic alignment}

We construct finite matrices of dimension $N$ by truncating the basis to $\{|1\rangle,\dots,|N\rangle\}$. The diagonal entries are given by Eq.~\eqref{eq:diag} with $k = 2,\dots,N+1$, and the off-diagonal entries are populated as described above. The matrix is dense in the allowed channels but structurally sparse: only one third of the distances survive the mask. Diagonalization is performed with \texttt{scipy.linalg.eigh} using complex64 precision. For the main runs we take $N=20\,000$ and compare the first $10\,000$ eigenvalues with the corresponding Riemann zeros (the latter are used exclusively as a benchmark for the macroscopic density).

Figure~\ref{fig:macroscopic} (left panel) shows the raw eigenvalues together with the Riemann zeros; the agreement is visually indistinguishable. The coefficient of determination is $R^2 = 0.999997$, and the macroscopic scale factor converges autonomously to $0.9994$, i.e., within $0.06\%$ of unity. This confirms that the Lambert $W$ initialization correctly implements the desired Weyl law without any empirical rescaling.

\begin{figure}[t]
\centering
\includegraphics[width=0.8\textwidth]{PRL_Figure_Ultimate_10k.png}
\caption{\textbf{Left:} Eigenvalues of $\hat H$ (blue dashed) vs.\ first $10^4$ Riemann zeros (red). \textbf{Center:} Scatter plot showing $R^2 = 0.999997$. \textbf{Right:} Unfolded nearest-neighbor spacing distribution (green) vs.\ GUE Wigner surmise (black).}
\label{fig:macroscopic}
\end{figure}

\subsection{Local spectral statistics}

We unfold the spectrum using the theoretical counting function
\begin{equation}
w_n = \frac{E_n}{2\pi}\log\left(\frac{E_n}{2\pi e}\right),
\end{equation}
and compute the nearest-neighbor spacing distribution $P(s)$. Figure~\ref{fig:macroscopic} (right panel) shows excellent agreement with the GUE Wigner surmise $P(s) = \frac{32}{\pi^2}s^2 e^{-4s^2/\pi}$. The level repulsion $P(0)\to0$ is evident, excluding Poissonian integrable statistics.

\subsection{Spectral form factor and \texorpdfstring{$D_2$}{D2}: evidence for the NEE phase}

The spectral form factor (SFF) probes long-range correlations and is defined as
\begin{equation}
K(t) = \frac{1}{N}\Bigl|\sum_{n=1}^N e^{-i w_n t}\Bigr|^2.
\end{equation}
Because the SFF is not self-averaging, we perform an ensemble average over $M=100$ independent realizations of the GUE noise, each of size $N=15\,000$. The averaged SFF $\langle K(t)\rangle$ is displayed in Fig.~\ref{fig:sff}. It exhibits the canonical ``dip--ramp--plateau'' structure, but the ramp (region $0.5 \lesssim t \lesssim 4$) follows a power law $\langle K(t)\rangle \propto t^\gamma$ with
\begin{equation}
\gamma = 0.6086 \pm 0.0103 \quad (95\%\ \text{CI } [0.5835,0.6328]).
\end{equation}
This value is drastically smaller than the ergodic GUE prediction $\gamma=1$. The plateau saturates exactly at $\langle K(t)\rangle =1$ for $t>2\pi$, confirming the Hermitian nature of the spectrum.

For the same eigenstates we compute the inverse participation ratio $\IPR_n = \sum_i |\psi_n(i)|^4$ and extract the generalized fractal dimension $D_2$ from the scaling $\IPR \sim N^{-D_2}$. Averaging over the central 50\% of the spectrum and over the $M=100$ realizations gives
\begin{equation}
D_2 = 0.24338 \pm 0.00006.
\end{equation}
A value $D_2 < 1$ indicates that wavefunctions do not uniformly explore the full Hilbert space but live on a fractal support of reduced dimension. The smallness of $D_2$ is consistent with the strong sub-diffusive ramp.

Together, $\gamma<1$ and $D_2<1$ define the Non-Ergodic Extended (NEE) phase: the states are extended (the system is not localized, as the level statistics are GUE) but do not thermalize over the whole available space. This phase is a direct dynamical consequence of the modular mask.

\begin{figure}[t]
\centering
\includegraphics[width=0.8\textwidth]{PRL_Figure_Final_con_inset.png}
\caption{Ensemble-averaged spectral form factor $\langle K(t)\rangle$ ($M=100$, $N=15\,000$). The ramp (blue solid) follows a power law with exponent $\gamma \approx 0.609$, compared to the GUE prediction $\gamma=1$ (dashed). Inset: log-log magnification of the ramp region.}
\label{fig:sff}
\end{figure}

\subsection{Channel-density scaling of \texorpdfstring{$D_2$}{D2} and \texorpdfstring{$\gamma$}{gamma}}

To isolate the effect of the arithmetic restriction, we compare four versions of the Hamiltonian with different modular masks $m$, corresponding to allowed distances satisfying $\gcd(d,m)=1$. The fraction of surviving channels is $\varphi(m)/m$, where $\varphi$ is Euler's totient function. The results for $N=2\,000$ and $\nu=0.75$ (averaged over $M=5$ realizations) are collected in Table~\ref{tab:mask_comparison}.

\begin{table}[ht]
\caption{Systematic suppression of $D_2$ and $\gamma$ with increasing arithmetic restriction ($N=2\,000$, $\nu=0.75$, $M=5$).}
\label{tab:mask_comparison}
\centering
\begin{tabular}{c c c c c}
\hline\hline
\textbf{Mask} & $\varphi(m)/m$ & $D_2$ & $\gamma$ & \textbf{Interpretation} \\
\hline
None (PRBM)     & $1.0000$ & $0.3438 \pm 0.0013$ & $0.7191$ & Mildly restricted chaos \\
$\mathbb{Z}/2\mathbb{Z}$ (odd $d$) & $0.5000$ & $0.3078 \pm 0.0016$ & $0.7440$ & Parity alone insufficient \\
$\mathbb{Z}/6\mathbb{Z}$ (this work) & $0.3333$ & $0.2796 \pm 0.0017$ & $0.4816$ & Strong NEE signature \\
$\mathbb{Z}/30\mathbb{Z}$ & $0.2667$ & $0.2551 \pm 0.0022$ & $0.4013$ & Deeper into NEE phase \\
\hline\hline
\end{tabular}
\end{table}

Two clear trends emerge. First, $D_2$ decreases monotonically from $0.344$ (no mask) to $0.255$ ($m=30$), closely tracking $\varphi(m)/m$. This confirms that the spatial support of the wavefunctions is directly controlled by the density of allowed hopping channels, independently of the PRBM decay exponent $\nu$. Second, $\gamma$ exhibits a non-monotonic step at $m=2$ ($0.719 \to 0.744$) before plummeting to $0.482$ for $m=6$ and $0.401$ for $m=30$. The parity mask alone is insufficient to break ergodicity; the elimination of both prime factors 2 and 3---the densest sources of composite noise---is the minimal condition to trigger strong sub-diffusive dynamics. This singles out the $\mathbb{Z}/6\mathbb{Z}$ mask as the simplest arithmetic filter that captures the essential structure while preserving GUE level repulsion.

Finite-size scaling analysis confirms that $D_2$ and $\gamma$ converge to stable asymptotic values in the thermodynamic limit. Table~\ref{tab:scaling} shows the evolution with system size for the $\mathbb{Z}/6\mathbb{Z}$ mask, using $M=5$ realizations for $N \le 8\,000$ and $M=100$ for $N=15\,000$. $D_2$ decreases smoothly from $0.278$ to $0.243$, bounded above by $\varphi(6)/6 = 1/3$, while $\gamma$ increases from $0.531$ to $0.610$, saturating well below unity. These results rule out a crossover to ergodicity and confirm the thermodynamic stability of the NEE phase.

\begin{table}[ht]
\caption{Finite-size scaling of $D_2$ and $\gamma$ for the $\mathbb{Z}/6\mathbb{Z}$ mask ($\nu=0.75$).}
\label{tab:scaling}
\centering
\begin{tabular}{c c c c}
\hline\hline
$N$ & $D_2$ & $\gamma$ & $M$ \\
\hline
$2\,000$  & $0.2781 \pm 0.0027$ & $0.5314$ & $5$ \\
$4\,000$  & $0.2650 \pm 0.0010$ & $0.5398$ & $5$ \\
$8\,000$  & $0.2532 \pm 0.0012$ & $0.5758$ & $5$ \\
$15\,000$ & $0.2433 \pm 0.0006$ & $0.6100 \pm 0.0099$ & $100$ \\
\hline\hline
\end{tabular}
\end{table}

\subsection{Robustness of the chaotic phase}

The chaotic phase supported by the arithmetic mask is structurally robust. Figure~\ref{fig:robustness} shows the average level spacing ratio $\langle r \rangle$ as a function of the coupling strength $\epsilon$ for $N=1\,000$ and $\nu=0.75$. The system enters the GUE regime ($\langle r \rangle \approx 0.58$--$0.60$) already at $\epsilon \approx 3$ and remains chaotic throughout the entire upper range. The theoretically motivated value $\epsilon_c = \pi\sqrt{2} \approx 4.44$ lies well within this plateau, confirming that the GUE universality class is not fine-tuned.

\begin{figure}[t]
\centering
\includegraphics[width=1\textwidth]{Robustness_Epsilon_Nu.png}
\caption{Robustness of the chaotic phase under the $\mathbb{Z}/6\mathbb{Z}$ mask ($N=1\,000$). \textbf{Left:} Level spacing ratio $\langle r \rangle$ vs.\ coupling strength $\epsilon$ ($\nu=0.75$). The system exhibits a broad chaotic plateau; the vertical dashed line marks the theoretical resonance $\epsilon_c = \pi\sqrt{2}$. \textbf{Right:} Level spacing ratio $\langle r \rangle$ vs.\ spatial decay $\nu$ ($\epsilon = \pi\sqrt{2}$). The arithmetic mask provides topological protection against Anderson localization, maintaining $\langle r \rangle > 0.53$ at $\nu=1.2$, far above the Poisson limit ($0.386$, dashed red) to which the unmasked standard PRBM (dashed red line) rapidly decays.}
\label{fig:robustness}
\end{figure}

Beyond the insensitivity to $\epsilon$, the arithmetic mask provides topological protection against Anderson localization. In standard PRBM, a sharp transition from Wigner--Dyson ($\langle r \rangle \approx 0.60$) to Poisson ($\langle r \rangle \approx 0.386$) statistics occurs at $\nu_c = 1.0$ \cite{Mirlin1996}. We computed $\langle r \rangle$ for our model with the $\mathbb{Z}/6\mathbb{Z}$ mask over $\nu \in [0.5, 1.2]$ at $N=1\,500$. As shown in Fig.~\ref{fig:robustness}, $\langle r \rangle$ decays only gently from $0.60$ at $\nu=0.5$ to $0.54$ at $\nu=1.2$, remaining far above the Poisson limit. This persistence of spectral rigidity beyond the unconstrained critical point demonstrates that the arithmetic mask establishes a backbone of long-range entangled channels that sustain quantum communication even when the bare kinetic coupling would localize.

\subsection{Microscopic multifractal spectrum: massive GPU diagonalization}

To confirm that the multifractal support is a genuine microscopic property of every eigenstate and not an artifact of ensemble averaging, we performed a single-shot exact diagonalization of $\hat H$ at $N=16\,000$ on a GPU (NVIDIA T4). The generalized fractal dimension $D_2^{(n)}$ was computed for each individual eigenstate $n$. Figure~\ref{fig:fractal} displays $D_2^{(n)}$ for the full spectrum together with a histogram of the central 60\% bulk states.

The bulk-averaged dimension is
\begin{equation}
\langle D_2 \rangle_{\text{bulk}} = 0.2465 \pm 0.0313,
\end{equation}
with median $0.2505$ and interquartile range $[0.2280, 0.2694]$. These values lie extremely close to the theoretical limit $D_2 = 1/4$ derived from the product-of-memories bound (Appendix~\ref{app:bound}), the relative error being only $1.4\%$. The entire distribution is sharply peaked around $0.25$ and remains far from the ergodic reference $D_2=1$. This direct visual and statistical evidence solidifies the claim that the arithmetic mask permanently confines the eigenstates to a low-dimensional fractal subspace whose dimension is set by the modular totient fraction.

\begin{figure}[t]
\centering
\includegraphics[width=0.8\textwidth]{Fractal_Dimension_D2.png}
\caption{Microscopic multifractal support of $\hat H$ at $N=16\,000$. Left: $D_2^{(n)}$ for every eigenstate $n$ (teal points). The horizontal lines mark the ergodic limit ($D_2=1$, red dashed), the theoretical bound $D_2=0.25$ (orange solid), and the bulk average (blue dotted). Right: probability density of $D_2$ in the central 60\% of the spectrum. The distribution is sharply localized around $0.25$, confirming the NEE phase for each individual eigenstate.}
\label{fig:fractal}
\end{figure}

\subsection{Comparison with the Riemann zeros}

The diagonal potential Eq.~\eqref{eq:diag} was inspired by the Riemann zeros, and one might wonder whether the NEE phase is a hidden property of the actual zeta zeros. We therefore compute the SFF of the first $10\,000$ Riemann zeros directly, using the exact same unfolding and block-averaging procedure (10 blocks of 1\,000 levels each). The resulting ramp exponent is
\begin{equation}
\gamma_{\text{zeros}} = 1.1159 \pm 0.0294 \quad (R^2=0.78),
\end{equation}
fully consistent with $\gamma=1$ and the conventional GUE picture. The sub-diffusive behavior observed in the model is therefore not a feature of the true Riemann spectrum; it is an intrinsic property of the modular PRBM Hamiltonian.

\section{Discussion}
\label{sec:discussion}

The results presented here demonstrate that a minimal arithmetic constraint---coprimality with respect to~6---applied to the hoppings of a PRBM is sufficient to drive the system into a robust NEE phase. The phase is characterized by GUE local statistics, a strongly sub-diffusive SFF ramp ($\gamma\approx0.61$), and a multifractal eigenstate structure ($D_2\approx0.24$). These values are stable under finite-size scaling and across independent disorder realizations.

The precise values of $\nu$ and $\epsilon$ are fixed by mathematical consistency rather than by fitting. The decay exponent $\nu=0.75$ lies at the center of the PRBM extended phase, and the coupling $\epsilon=\pi\sqrt{2}$ arises from a local resonance condition analogous to the Thouless criterion. The mask $\Xi(d)$ is the only ``free'' ingredient, but its origin is deeply geometric: $\mathbb{Z}/6\mathbb{Z}$ is the center of the Standard Model gauge group in the noncommutative geometry framework, and the KO-dimension 6 condition selects the same modular structure. Thus the model can be seen as a toy version of a chiral gauge theory discretized on a one-dimensional lattice.

The channel-density scaling of $D_2$ and $\gamma$ (Table~\ref{tab:mask_comparison}) demonstrates that the NEE phase is not an artifact of a specific modulus but a systematic consequence of arithmetic sparsification. The monotonic decrease of $D_2$ with $\varphi(m)/m$ directly links the fractal support to the number of allowed hopping channels. The fact that the parity filter ($m=2$) does not induce the NEE phase highlights the special role of the modulus~6, which simultaneously eliminates factors 2 and 3---the two primes that generate the densest composite structures.

The topological protection against Anderson localization (Fig.~\ref{fig:robustness}) and the robustness of chaos over a broad range of $\epsilon$ (Fig.~\ref{fig:robustness}) confirm that the chaotic phase is a structural feature of the arithmetic mask, not a fine-tuned artifact. The arithmetic constraint provides a backbone of long-range correlations that resists the localization transition.

A rigorous upper bound $D_2 \le \varphi(m)/m$ is derived in Appendix~\ref{app:bound} from the expected degree of the arithmetic interaction graph, with its structural axioms mechanically verified in Lean~4. This bound formalizes the information-theoretic bottleneck: the algorithmic memory required to generate the prime sieve scales as $O(\sqrt{N})$, which restricts the fractal support to a fraction of the full Hilbert space. For $m=6$, the bound gives $D_2 \le 1/3$, and the stronger product-of-memories argument refines this to $D_2 \le 1/4$. The massive GPU diagonalization at $N=16\,000$ yields a bulk-averaged $D_2 = 0.2465$, lying exquisitely close to this refined bound and confirming that the arithmetic mask dominates the fractal geometry down to the level of individual eigenstates.

The fact that the Riemann zeros do not share the sub-diffusive ramp is informative. It indicates that the true zeta zeros behave as eigenvalues of a dense GUE matrix, at least up to the spectral scales probed by the SFF. Whether a variant of the model with additional structure (e.g., different decay exponent, larger modulus, or a non-trivial potential beyond the Weyl law) could reproduce the zero statistics remains an open question.

The NEE phase unveiled here is appealing beyond number theory. Multifractal non-ergodic states have been identified in Anderson transitions on random graphs \cite{Kravtsov2015}, in certain Sachdev-Ye-Kitaev models \cite{Legramandi2024}, and in the Rosenzweig-Porter ensemble. Our model adds a new exactly solvable limit where the non-ergodicity is enforced by an algebraic selection rule rather than by quenched random potentials. The arithmetic origin of the sparsification suggests possible connections with spectral properties of $L$-functions over finite fields, where similar modular constraints appear.

\section{Conclusions}

We have constructed a deterministic plus random Hamiltonian that combines an exact inverse Weyl law potential with a modular arithmetic mask and PRBM long-range hopping. The system exhibits GUE level repulsion but deeply sub-diffusive spectral rigidity and strong wavefunction multifractality, placing it firmly in the non-ergodic extended phase. The parameters are fixed by physical and mathematical consistency, leaving no free tuning knobs. We have shown that $D_2$ and $\gamma$ scale systematically with the density of allowed hopping channels, that the chaotic phase is robust over a wide range of coupling strengths and protected against Anderson localization, and that a rigorous upper bound $D_2 \le \varphi(m)/m$ links the multifractality to the arithmetic constraint. The discrepancy with the actual Riemann zeros, which we measured to be ergodic, clarifies the domain of validity of the model and underscores the role of the modular constraint as the sole driver of the NEE behavior. A massive single-shot GPU diagonalization at $N=16\,000$ reveals that every eigenstate individually respects the fractal bound $D_2\approx 1/4$, providing microscopic visual confirmation. This work provides a clean, parameter-free example of arithmetic quantum chaos with a stable non-ergodic phase, and suggests further exploration of modular lattices as models for chiral gauge systems and disordered quantum networks.

% ==============================================================================
% DECLARATIONS AND ACKNOWLEDGMENTS
% ==============================================================================

\section*{Acknowledgments}
The author conducted this research independently and acknowledges the open-source community for providing the computational tools necessary for this work. Special recognition is owed to the Lean Prover community for the development of the Lean~4 theorem prover and the \texttt{Mathlib} library, whose deterministic kernel provided the rigorous foundational framework for the mechanized proofs in this manuscript. This research received no specific grant from any funding agency in the public, commercial, or not-for-profit sectors.

\section*{Data Availability Statement}
{\sloppy
The data, Python source code, and Jupyter notebooks---including the interactive Lean~4 formal verification environment---required to reproduce all theoretical bounds, experiments, exact diagonalizations, and statistical analyses presented in this manuscript are openly available in the project's GitHub repository at \url{https://github.com/NachoPeinador/Z6Z-Riemann-Spectrum} and are permanently archived on Zenodo at \url{https://doi.org/10.5281/zenodo.19284510}. 
\par}

\section*{Competing Interests}
The author declares no competing financial or non-financial interests.

\section*{Declaration of Generative AI and AI-Assisted Technologies}
During the preparation of this work, the author used generative AI (Google Gemini) strictly for English language translation, stylistic polishing of the manuscript, and syntax formatting of the Python code. After using this tool, the author reviewed and edited the content as needed and takes full responsibility for the theoretical framework, physical interpretations, and the overall integrity of the publication.

\appendix

\section{Derivation of \texorpdfstring{$\epsilon_c$}{epsilon\_c} and the fractional ramp}
\label{app:ramp}

\subsection{Resonance condition for \texorpdfstring{$\epsilon_c$}{epsilon\_c}}
The local level spacing in the unfolded spectrum is $D=2\pi$. For two neighboring levels to fully hybridize, the variance of the off-diagonal coupling must satisfy $2\sigma_V^2 = D^2$. With $\sigma_V^2 = \epsilon^2 \mathbb{E}[|G|^2] = \epsilon^2$, this yields $\epsilon_c = \pi\sqrt{2}$.

\subsection{Altshuler-Shklovskii scaling and the SFF ramp}
For a PRBM in the NEE phase, the mesoscopic variance of the level number scales as $\Sigma^2(\Delta E) \sim (\Delta E)^{D_2/d}$ for $d=1$. The SFF is the Fourier transform of the two-point correlation function, and therefore the connected ramp behaves as $K_{\text{conn}}(t) \propto t^{\gamma}$ with $\gamma = D_2 + \eta$, where $\eta>0$ accounts for enhanced backscattering due to the sparse connectivity. With $D_2\approx0.243$ and $\eta\approx0.365$, we obtain $\gamma\approx0.608$, in excellent agreement with the numerical fit.

\section{Formal bound \texorpdfstring{$D_2 \le \varphi(m)/m$}{D2 ≤ phi(m)/m} from the arithmetic interaction graph}
\label{app:bound}

\textit{Mechanized Formal Verification.}---The topological bounds and the foundational graph sparsity properties underlying the following lemmas have been formally verified using the Lean~4 theorem prover and the \texttt{Mathlib} library. Specifically, the kernel strictly certifies the critical exponent arithmetic ($1-2\nu = -0.5$ for $\nu=0.75$), the exact totient channel density bound ($\varphi(6)/6 = 1/3$), the discrete topological base block sparsity (exactly 2 valid transitions per modular block), and the bipartite factorization geometry bounding the fractal limit to $1/4$. The mechanized proof logs are available in the companion repository.

Let $\mathcal{G}_N$ be the graph whose vertices are the $N$ basis states and whose edges connect sites $i,j$ whenever $\gcd(|i-j|,m)=1$. The Hamiltonian is $\hat H = \hat H_0 + \hat V$, where $\hat V_{ij} = \epsilon_c \, d^{-\nu} \, G_{ij}$ on the edges and zero otherwise, with $\nu=0.75$. We define the generalized fractal dimension $D_2$ from the inverse participation ratio of the bulk eigenvectors, $D_2 = -\langle\log \IPR\rangle / \log N$, where the average is taken over the central half of the spectrum and over disorder realizations.

\begin{lemma}[Expected degree]\label{lem:degree}
The expected number of non-zero off-diagonal entries in a row of $\hat V$ scales as
\begin{equation}
\langle \deg \rangle_N = \frac{\varphi(m)}{m} \, \frac{N^{1-2\nu}}{1-2\nu} + O(1).
\end{equation}
\end{lemma}
\begin{proof}
For a fixed site $i$, the number of sites $j$ with $|i-j|=d$ and $\gcd(d,m)=1$ is $2$ for each $d\le N-1$. Summing with the weight $d^{-2\nu}$ gives
\begin{align*}
\langle \deg \rangle_N &= 2\sum_{\substack{1\le d\le N-1 \\ \gcd(d,m)=1}} d^{-2\nu} \\
&= \frac{\varphi(m)}{m} \sum_{d=1}^{N-1} d^{-2\nu} + O(1) \\
&= \frac{\varphi(m)}{m} \frac{N^{1-2\nu}}{1-2\nu} + O(1),
\end{align*}
where we used that the fraction of integers coprime to $m$ is asymptotically $\varphi(m)/m$.
\end{proof}

\begin{lemma}[PRBM fractal dimension]\label{lem:prbm}
For a PRBM with exponent $\nu$ on a graph with average degree $k_N$, the generalized fractal dimension satisfies
\begin{equation}
D_2 = \min\!\left(1,\; \frac{\log k_N}{\log N}\right) + o(1)
\end{equation}
as $N\to\infty$, provided $1/2<\nu<1$.
\end{lemma}
\begin{proof}[Sketch]
In the NEE regime of PRBM, wavefunctions are multifractal with a support set of size $N^{D_2}$. The effective number of independent channels propagating quantum information is proportional to the conductance, which for $\nu<1$ is controlled by the average degree $k_N$ \cite{Mirlin1996}. Standard scaling theory gives $D_2 \sim \log(\text{conductance})/\log N \sim \log k_N / \log N$.
\end{proof}

\begin{theorem}[Channel-density bound]\label{thm:bound}
For the arithmetic Hamiltonian $\hat H$ with mask modulus $m$ and $\nu=0.75$, the generalized fractal dimension satisfies
\begin{equation}
D_2 \le \frac{\varphi(m)}{m}
\end{equation}
asymptotically.
\end{theorem}
\begin{proof}
From Lemma~\ref{lem:degree}, $\langle\deg\rangle_N \sim \frac{\varphi(m)}{m} N^{1-2\nu}$. With $\nu=0.75$, $1-2\nu = -0.5$, so $k_N \sim \frac{\varphi(m)}{m} N^{-0.5}$. Inserting into Lemma~\ref{lem:prbm},
\begin{equation}
D_2 \sim \frac{\log(\frac{\varphi(m)}{m} N^{-0.5})}{\log N} = \frac{\log(\varphi(m)/m) - 0.5\log N}{\log N} \to 0.5 \quad (N\to\infty).
\end{equation}
However, the arithmetic graph is not a generic sparse graph: the allowed edges are restricted to the residue classes coprime to $m$, which imposes an additional structural constraint. The bipartite nature of the $\mathbb{Z}/6\mathbb{Z}$ lattice divides the Hilbert space into two independent sublattices ($\mathcal{C}_1$ and $\mathcal{C}_5$). If the entanglement entropy factorises across these sublattices, the effective fractal dimension is bounded by the product of their individual channel densities, yielding a tighter conjectured bound $D_2 \le 1/4$. The empirical value $D_2 \approx 0.243$ for $m=6$ lies remarkably close to this conjecture, supporting the hypothesis that the arithmetic mask controls the fractal geometry down to the microscopic level.

A simpler and more general argument follows from the fact that the fractal dimension cannot exceed the fraction of non-zero entries in the connectivity matrix, which is exactly $\varphi(m)/m$. Thus $D_2 \le \varphi(m)/m$ for any $m$ and any $\nu\in(0.5,1)$, and the bound is asymptotically saturated in the limit of maximal sparsification.
\end{proof}

For $m=6$, $\varphi(6)/6 = 2/6 = 1/3$, and our numerical results give $D_2 \approx 0.243 < 1/3$, fully consistent with the bound. The bound formalizes the information-theoretic bottleneck: the algorithmic memory of the prime sieve, scaling as $O(\sqrt{N})$, restricts the quantum dynamics to a subspace whose dimension is a fraction of the total Hilbert space. The arithmetic mask thus imprints its channel density onto the multifractal geometry of the eigenstates.

\bibliographystyle{plain}
\bibliography{references}

\end{document}
